% !TeX root = ../../Book1.tex
% 实读来源：Fremlin, Measure Theory, Vol. 2，作者在线第28章摘本，
% 285A、285F、285I、285J、285L--M；摘本 PDF 第67、69--73页。
% 条号作为正式定位，不把持续修订摘本的 PDF 页码当作2016版页码。
% Bass, Real Analysis for Graduate Students, v5，印刷第297--302、305--307页，
% 仅作教学背景。以下 Gaussian 平滑唯一性与极限存在性均在正文自证。
\section{特征函数与概率弱收敛}
\label{b1:sec:3006}

大数律比较随机变量与一个常数的距离；中心极限定理则要比较整个分布。
直接对所有有界连续函数计算积分，通常不便于处理独立和。
复指数函数的优势在于：它始终有界，并且把加法变成乘法。
本节先用这些特殊的测试函数表示分布，再说明它们为什么足以辨认分布及其弱极限。

\subsection{特征函数与 Gaussian 模型}
\label{b1:subsec:300601}

对实数 \(u\)，记 \(e^{iu}=\cos u+i\sin u\)。
这里的复值积分按实部、虚部分别定义；由于 \(|e^{iu}|=1\)，
对任意概率分布都可以积分，不需要假设一阶矩存在。

\begin{definition}\label{b1:def:30-characteristic-function}
设 \(\mu\) 是 \(\mathbb R\) 上的 Borel 概率测度。
函数
\[
 \varphi_\mu(t)=\int_{\mathbb R}e^{itx}\,\mathrm d\mu(x),
 \quad t\in\mathbb R,
\]
称为 \(\mu\) 的\term[tezheng-hanshu]{特征函数}[characteristic function]。
若 \(X\) 是概率空间 \((\Omega,\mathcal F,\mathbb P)\) 上的实随机变量，
其分布为 \(\mu_X\)，则规定
\[
 \varphi_X(t)=\varphi_{\mu_X}(t)=\mathbb E(e^{itX}).
\]
\end{definition}
\symidx{\(\varphi_\mu,\varphi_X\)：概率分布或随机变量的特征函数}

这一约定在指数中取正号，且不加归一化系数。
特征函数的定义和下面的基本运算可对照
Fremlin 的《Measure Theory》\cite[285A, 285F, 285I]{Fremlin2016Measure}。

\begin{proposition}\label{b1:prop:30-characteristic-basic}
设 \(\mu\) 是实轴上的 Borel 概率测度，则
\[
 \varphi_\mu(0)=1,\quad
 |\varphi_\mu(t)|\leq1,\quad
 \varphi_\mu(-t)=\overline{\varphi_\mu(t)}
 \quad(t\in\mathbb R),
\]
且 \(\varphi_\mu\) 在 \(\mathbb R\) 上一致连续。
若 \(X\) 是任意概率空间上的实随机变量，\(a,b\in\mathbb R\)，则
\[
 \varphi_{aX+b}(t)=e^{itb}\cdot\varphi_X(at).
\]
若 \(X_1,\ldots,X_r\) 是同一概率空间上的独立实随机变量，
其中整数 \(r\geq1\)，则
\begin{equation}\label{b1:eq:30-independent-characteristic}
 \varphi_{X_1+\cdots+X_r}(t)
 =\prod_{j=1}^r\varphi_{X_j}(t)
 \quad(t\in\mathbb R).
\end{equation}
\end{proposition}
\begin{proof}
前三式分别由总质量为一、积分三角不等式和共轭与积分交换得到。
对任意 \(t,h\in\mathbb R\)，有
\[
 |\varphi_\mu(t+h)-\varphi_\mu(t)|
 \leq\int_{\mathbb R}|e^{ihx}-1|\,\mathrm d\mu(x).
\]
当 \(h\to0\) 时，被积函数逐点趋于零，并由可积常数 \(2\) 控制。
控制收敛定理使右端趋于零；右端又与 \(t\) 无关，故得到一致连续性。

仿射变换公式来自
\(e^{it(aX+b)}=e^{itb}\cdot e^{i(at)X}\)。
对独立变量，\cref{b1:prop:30-independent-joint-law}给出联合分布
\(\mu_{X_1}\otimes\cdots\otimes\mu_{X_r}\)。
被积函数 \(\prod_{j=1}^r e^{itx_j}\) 的模恒为一，
因此在这个乘积概率空间上绝对可积。
Fubini 定理遂将其积分拆成各个特征函数的乘积，
得到\eqref{b1:eq:30-independent-characteristic}。
\end{proof}

单点分布 \(\delta_a\) 的特征函数是 \(e^{ita}\)；
参数 \(p\in[0,1]\) 的 Bernoulli 分布的特征函数是
\(1-p+p e^{it}\)。
接下来需要一个连续分布模型，它的特征函数仍具有便于取幂的形式。
Gaussian 密度具有这一性质，下面计算它的特征函数。

\begin{lemma}\label{b1:lem:30-gaussian-fourier}
对每个 \(\varepsilon>0\)，定义
\[
 g_\varepsilon(x)=\frac{1}{\sqrt{2\pi}\,\varepsilon}
       \exp\left(-\frac{x^2}{2\varepsilon^2}\right),
 \quad x\in\mathbb R.
\]
则 \(g_\varepsilon\) 是严格正的连续概率密度，并满足
\[
 \int_{\mathbb R}x\cdot g_\varepsilon(x)\,\mathrm dx=0,
 \quad
 \int_{\mathbb R}x^2\cdot g_\varepsilon(x)\,\mathrm dx=\varepsilon^2.
\]
此外，对所有 \(t,x\in\mathbb R\)，有
\begin{align}
 \int_{\mathbb R}e^{itx}\cdot g_\varepsilon(x)\,\mathrm dx
   &=e^{-\varepsilon^2t^2/2},
       \label{b1:eq:30-gaussian-characteristic}\\
 g_\varepsilon(x)
   &=\frac{1}{2\pi}\int_{\mathbb R}
       e^{-itx}\cdot e^{-\varepsilon^2t^2/2}\,\mathrm dt.
       \label{b1:eq:30-gaussian-inversion}
\end{align}
上述各积分均绝对收敛。
\end{lemma}
\begin{proof}
\secref{b1:sec:0705}已用 Tonelli 定理和极坐标换元证明
\(\int_{\mathbb R}e^{-u^2}\,\mathrm du=\sqrt\pi\)。
在此公式中作线性换元 \(x=\sqrt2\,\varepsilon u\)，得到
\[
 \int_{\mathbb R}e^{-x^2/(2\varepsilon^2)}\,\mathrm dx
 =\sqrt2\,\varepsilon\int_{\mathbb R}e^{-u^2}\,\mathrm du
 =\sqrt{2\pi}\,\varepsilon.
\]
这说明 \(g_\varepsilon\) 可积且积分为一；其正性、连续性由公式直接给出。

多项式乘 Gaussian 的绝对可积性来自指数衰减；
例如在充分大的 \(|x|\) 上，
\(|x|^k e^{-x^2/2}\leq C_k e^{-|x|/4}\)，其中 \(k=1,2\)。
奇偶性给出 \(\int x\cdot g_1(x)\,\mathrm dx=0\)。
又有 \(g_1'(x)=-x\cdot g_1(x)\)，故在 \([-R,R]\) 上分部积分，令 \(R\to\infty\)，得到
\[
 \int_{\mathbb R}x^2\cdot g_1(x)\,\mathrm dx
 =-\bigl[x\cdot g_1(x)\bigr]_{-\infty}^{+\infty}
     +\int_{\mathbb R}g_1(x)\,\mathrm dx
 =1.
\]
这里边界项趋于零，各积分绝对收敛。
换元 \(x=\varepsilon u\) 即给出一般 \(\varepsilon\) 的两个矩公式。

令
\[
 I(t)=\int_{\mathbb R}e^{itx}\cdot g_1(x)\,\mathrm dx.
\]
实参数差商满足
\[
 \left|\frac{e^{i(t+h)x}-e^{itx}}{h}\right|\leq|x|
 \quad(h\ne0).
\]
由于 \(|x|g_1(x)\) 可积，控制收敛允许对实部、虚部在积分号下求导：
\[
 I'(t)=i\int_{\mathbb R}x\cdot e^{itx}\cdot g_1(x)\,\mathrm dx.
\]
再在 \([-R,R]\) 上分部积分，并利用 \(g_1'=-x\cdot g_1\)，得到
\[
 \int_{\mathbb R}x\cdot e^{itx}\cdot g_1(x)\,\mathrm dx
 =-\bigl[e^{itx}\cdot g_1(x)\bigr]_{-\infty}^{+\infty}
       +it\int_{\mathbb R}e^{itx}\cdot g_1(x)\,\mathrm dx
 =it\cdot I(t).
\]
因而 \(I'(t)=-tI(t)\)，且 \(I(0)=1\)。
函数 \(e^{t^2/2}\cdot I(t)\) 的导数为零，于是
\(I(t)=e^{-t^2/2}\)。
用 \(x=\varepsilon u\) 换元便得到\eqref{b1:eq:30-gaussian-characteristic}。

最后，以 \(s=\varepsilon t\) 换元，再次使用刚求出的 \(I\)，可得
\[
 \int_{\mathbb R}e^{-itx}\cdot e^{-\varepsilon^2t^2/2}\,\mathrm dt
 =\frac{\sqrt{2\pi}}{\varepsilon}\cdot I(-x/\varepsilon)
 =\frac{\sqrt{2\pi}}{\varepsilon}
       \exp\left(-\frac{x^2}{2\varepsilon^2}\right).
\]
除以 \(2\pi\)，即是\eqref{b1:eq:30-gaussian-inversion}。
其中指数因子的模为一，剩余 Gaussian 因子可积，
所以最后的积分同样绝对收敛。
\end{proof}

\begin{definition}\label{b1:def:30-normal-distribution}
设 \(m\in\mathbb R\)、\(\sigma\geq0\)。
当 \(\sigma>0\) 时，把具有密度
\[
 x\longmapsto g_\sigma(x-m)
 =\frac{1}{\sqrt{2\pi}\,\sigma}
       \exp\left(-\frac{(x-m)^2}{2\sigma^2}\right)
\]
的 Borel 概率分布称为均值 \(m\)、方差 \(\sigma^2\) 的
\term[zhengtai-fenbu]{正态分布}[normal distribution]，
记为 \(N(m,\sigma^2)\)。
约定 \(N(m,0)=\delta_m\)。
其中 \(N(0,1)\) 称为
\term[biaozhun-zhengtai-fenbu]{标准正态分布}[standard normal distribution]，
其分布函数记为
\[
 \Phi(x)=\frac{1}{\sqrt{2\pi}}\int_{-\infty}^x e^{-u^2/2}\,\mathrm du.
\]
\end{definition}
\symidx{\(N(m,\sigma^2)\)：均值 \(m\)、方差 \(\sigma^2\) 的正态分布}
\symidx{\(\Phi\)：标准正态分布函数}

若 \(G\) 的分布为 \(N(0,1)\)，则 \(m+\sigma G\) 的分布为
\(N(m,\sigma^2)\)；当 \(\sigma>0\) 时由密度换元得到，
当 \(\sigma=0\) 时直接得到单点分布。
Gaussian 引理同时验证了名称中的均值与方差，并给出
\begin{equation}\label{b1:eq:30-normal-characteristic}
 \varphi_{N(m,\sigma^2)}(t)
 =\exp\left(imt-\frac{\sigma^2t^2}{2}\right).
\end{equation}
\(\Phi\) 连续且严格递增，满足
\(\Phi(-x)=1-\Phi(x)\)，在两端的极限分别为 \(0\)、\(1\)。
这些性质都来自正密度 \(g_1\) 的连续性、对称性及积分为一。

\subsection{特征函数的唯一性}
\label{b1:subsec:300602}

两个分布若对所有复指数函数给出相同积分，为什么必定相等？
我们先给分布加上一层 Gaussian 平滑。
平滑后的密度可以由特征函数的一个绝对收敛积分写出；
随后让平滑尺度趋于零，就能恢复原来的分布。

\begin{lemma}\label{b1:lem:30-gaussian-smoothing}
设 \(\mu\) 是实轴上的 Borel 概率测度，\(\varepsilon>0\)，并令
\[
 p_{\mu,\varepsilon}(x)
   =\int_{\mathbb R}g_\varepsilon(x-y)\,\mathrm d\mu(y),
 \quad x\in\mathbb R.
\]
则 \(p_{\mu,\varepsilon}\) 是连续概率密度，且
\begin{equation}\label{b1:eq:30-smoothed-density}
 p_{\mu,\varepsilon}(x)
 =\frac{1}{2\pi}\int_{\mathbb R}
       e^{-itx}\cdot\varphi_\mu(t)\cdot e^{-\varepsilon^2t^2/2}\,\mathrm dt.
\end{equation}
右端绝对收敛。若
\(\rho_{\mu,\varepsilon}=p_{\mu,\varepsilon}(x)\,\mathrm dx\)，则
\(\rho_{\mu,\varepsilon}\Rightarrow\mu\) 当 \(\varepsilon\downarrow0\)。
\end{lemma}
\begin{proof}
被积函数非负，Tonelli 定理与 \(\int g_\varepsilon=1\) 给出
\[
 \int_{\mathbb R}p_{\mu,\varepsilon}(x)\,\mathrm dx
 =\int_{\mathbb R}\left(\int_{\mathbb R}
       g_\varepsilon(x-y)\,\mathrm dx\right)\mathrm d\mu(y)
 =1.
\]
当 \(x_n\to x\) 时，\(g_\varepsilon(x_n-y)\to g_\varepsilon(x-y)\)，
而这些函数都由常数 \((\sqrt{2\pi}\,\varepsilon)^{-1}\) 控制。
该常数对概率测度 \(\mu\) 可积，故控制收敛给出密度的连续性。

将\eqref{b1:eq:30-gaussian-inversion}代入定义，可写成
\[
 p_{\mu,\varepsilon}(x)
 =\frac{1}{2\pi}\int_{\mathbb R}\int_{\mathbb R}
       e^{-it(x-y)}\cdot e^{-\varepsilon^2t^2/2}
       \,\mathrm dt\,\mathrm d\mu(y).
\]
交换积分所需的绝对积分为
\[
 \int_{\mathbb R}\int_{\mathbb R}
       e^{-\varepsilon^2t^2/2}\,\mathrm dt\,\mathrm d\mu(y)
 =\int_{\mathbb R}e^{-\varepsilon^2t^2/2}\,\mathrm dt<\infty.
\]
因此 Fubini 定理适用，先对 \(y\) 积分便得到
\eqref{b1:eq:30-smoothed-density}。

为证明弱收敛，任取实值 \(f\in C_b(\mathbb R)\)。
Fubini 定理及换元 \(x=y+\varepsilon z\) 给出
\[
 \int_{\mathbb R}f\,\mathrm d\rho_{\mu,\varepsilon}
 =\int_{\mathbb R}\int_{\mathbb R}
       f(y+\varepsilon z)\cdot g_1(z)\,\mathrm dz\,\mathrm d\mu(y).
\]
可以把右端看作固定乘积概率测度
\(\mu(\mathrm dy)\otimes(g_1(z)\,\mathrm dz)\) 上的积分。
对每个 \(y,z\)，连续性给出 \(f(y+\varepsilon z)\to f(y)\)，
而其绝对值不超过 \(\|f\|_\infty\)。
控制收敛遂给出极限 \(\int f\,\mathrm d\mu\)。
这恰好是第12章的弱收敛定义；这里不要求 \(f\) 在整个实轴上一致连续。
\end{proof}

\begin{theorem}[特征函数的唯一性]\label{b1:thm:30-characteristic-uniqueness}
设 \(\mu,\nu\) 是 \(\mathbb R\) 上的 Borel 概率测度。
若对每个 \(t\in\mathbb R\) 都有 \(\varphi_\mu(t)=\varphi_\nu(t)\)，
则 \(\mu=\nu\)。
\end{theorem}
\begin{proof}
由\eqref{b1:eq:30-smoothed-density}，对每个 \(\varepsilon>0\)，
都有 \(p_{\mu,\varepsilon}=p_{\nu,\varepsilon}\)，从而
\(\rho_{\mu,\varepsilon}=\rho_{\nu,\varepsilon}\)。
对任意实值 \(f\in C_b(\mathbb R)\)，使用平滑引理令 \(\varepsilon\downarrow0\)，得到
\[
 \int_{\mathbb R}f\,\mathrm d\mu
 =\lim_{\varepsilon\downarrow0}\int_{\mathbb R}f\,\mathrm d\rho_{\mu,\varepsilon}
 =\lim_{\varepsilon\downarrow0}\int_{\mathbb R}f\,\mathrm d\rho_{\nu,\varepsilon}
 =\int_{\mathbb R}f\,\mathrm d\nu.
\]
实轴是度量空间，\(\mu,\nu\) 都是有限正 Borel 测度，
所以\cref{b1:prop:weak-mass-uniqueness}给出 \(\mu=\nu\)。
\end{proof}

因此，特征函数并未丢失分布信息。
例如，要证明一个独立和仍服从某个正态分布，
只需计算其特征函数，并与\eqref{b1:eq:30-normal-characteristic}比较。
但在处理一列分布时，除了辨认候选极限，还必须保证质量没有逃向无穷远。

\subsection{紧性与子列极限的辨认}
\label{b1:subsec:300603}

第12章的紧性要求同一个紧集控制全族的尾质量。
特征函数如何提供这样的控制？在原点附近，
\(1-\operatorname{Re}\varphi_\mu(t)\) 是 \(1-\cos(tx)\) 的平均值。
再对一个短参数区间求平均，便能避免单个余弦函数不断回到一的困难。
这一尾概率估计的同类形式见
Fremlin\cite[285J]{Fremlin2016Measure}；下面给出本节所用的常数和证明。

\begin{lemma}\label{b1:lem:30-characteristic-tail}
设 \(\mu\) 是实轴上的 Borel 概率测度，\(\delta>0\)。
则
\begin{equation}\label{b1:eq:30-characteristic-tail}
 \mu\bigl(\{x:|x|\geq2/\delta\}\bigr)
 \leq\frac{1}{\delta}\int_{-\delta}^{\delta}
       \bigl(1-\operatorname{Re}\varphi_\mu(t)\bigr)\,\mathrm dt.
\end{equation}
\end{lemma}
\begin{proof}
函数 \(1-\cos(tx)\) 非负且不超过 \(2\)，
在 \([-\delta,\delta]\times\mathbb R\) 上对相应乘积测度可积。
交换积分并计算内层积分，得到
\[
 \frac{1}{2\delta}\int_{-\delta}^{\delta}
       \bigl(1-\operatorname{Re}\varphi_\mu(t)\bigr)\,\mathrm dt
 =\int_{\mathbb R}\left(1-\frac{\sin(\delta x)}{\delta x}\right)\mathrm d\mu(x),
\]
其中在 \(x=0\) 处将比值取为连续值 \(1\)。
右端被积函数处处非负，因为它也是 \(1-\cos(tx)\) 的区间平均。
当 \(|x|\geq2/\delta\) 时，
\(\sin(\delta x)/(\delta x)\leq1/2\)，故被积函数至少为 \(1/2\)。
于是右端至少为
\(\frac12\cdot\mu(\{x:|x|\geq2/\delta\})\)，这就给出所需不等式。
\end{proof}

\begin{lemma}\label{b1:lem:30-characteristic-tightness}
设 \((\mu_n)_{n\geq1}\) 是实轴上的 Borel 概率测度列。
若对每个 \(t\in\mathbb R\)，
\(\varphi_{\mu_n}(t)\to\psi(t)\)，且函数 \(\psi:\mathbb R\to\mathbb C\)
在 \(0\) 处连续，则 \(\{\mu_n:n\geq1\}\) 是紧族。
\end{lemma}
\begin{proof}
因为各 \(\varphi_{\mu_n}(0)=1\)，所以 \(\psi(0)=1\)。
给定 \(\eta>0\)，由零点处连续性，取 \(\delta>0\)，使
\[
 |1-\psi(t)|<\eta/8\quad(|t|\leq\delta).
\]
极限函数可测，且 \(|\psi(t)|\leq1\)。
在有限区间 \([-\delta,\delta]\) 上，
\(0\leq1-\operatorname{Re}\varphi_{\mu_n}(t)\leq2\)，
故控制收敛给出
\[
 \frac{1}{\delta}\int_{-\delta}^{\delta}
       \bigl(1-\operatorname{Re}\varphi_{\mu_n}(t)\bigr)\,\mathrm dt
 \longrightarrow
 \frac{1}{\delta}\int_{-\delta}^{\delta}
       \bigl(1-\operatorname{Re}\psi(t)\bigr)\,\mathrm dt
 \leq\eta/4.
\]
因此存在 \(N\)，使所有 \(n\geq N\) 的左端都小于 \(\eta/2\)。
由尾概率估计，这些 \(\mu_n\) 在
\(\{|x|\geq2/\delta\}\) 上的质量都小于 \(\eta/2\)。

还须把前面有限个测度纳入同一个控制。
对每个 \(n<N\)，由
\([-R,R]\uparrow\mathbb R\) 及测度连续性，可取足够大的 \(R_n\)，
使 \(\mu_n(\mathbb R\setminus[-R_n,R_n])<\eta\)。
取 \(R\) 不小于 \(2/\delta\) 及这些有限个 \(R_n\)，便有
\[
 \sup_{n\geq1}\mu_n(\mathbb R\setminus[-R,R])<\eta.
\]
闭区间 \([-R,R]\) 紧，所以满足\cref{b1:def:tight-family}的要求。
\end{proof}

\begin{proposition}\label{b1:prop:30-tight-characteristic-limit}
设 \((\mu_n)_{n\geq1}\) 是实轴上的 Borel 概率测度列，
\(\mu\) 是实轴上的 Borel 概率测度。
若 \(\{\mu_n:n\geq1\}\) 为紧族，且
\(\varphi_{\mu_n}(t)\to\varphi_\mu(t)\) 对每个 \(t\in\mathbb R\) 成立，
则 \(\mu_n\Rightarrow\mu\)。
\end{proposition}
\begin{proof}
实轴在通常度量下完备且可分，因而是波兰空间。
由\cref{b1:thm:prokhorov}，任何子列都有一个弱收敛的子子列；
设其中一个子子列的极限为 Borel 概率测度 \(\nu\)。
对每个固定 \(t\)，函数 \(\cos(tx)\)、\(\sin(tx)\) 都属于
\(C_b(\mathbb R)\)，所以弱收敛给出该子子列的特征函数趋于
\(\varphi_\nu(t)\)。
原来的点态极限假设又给出 \(\varphi_\nu(t)=\varphi_\mu(t)\)。
唯一性定理于是说明 \(\nu=\mu\)。
所有弱收敛子列的极限都被辨认为同一个 \(\mu\)，
再由\cref{b1:cor:tight-unique-limit}得到整列弱收敛。
\end{proof}

\subsection{特征函数的连续性定理}
\label{b1:subsec:300604}

前一命题已经假定候选概率测度 \(\mu\) 存在。
Fremlin\cite[285L--M]{Fremlin2016Measure}讨论了给定极限测度后的弱收敛判别和唯一性。
对中心极限定理而言，我们最先算出的却只是一个函数极限。
下面的\term[levy-lianxuxing-dingli]{Lévy 连续性定理}[Lévy's continuity theorem]
同时解决极限分布的存在与辨认。

\begin{theorem}[Lévy 连续性定理]\label{b1:thm:30-levy-continuity}
设 \((\mu_n)_{n\geq1}\) 是 \(\mathbb R\) 上的 Borel 概率测度列。
若存在函数 \(\psi:\mathbb R\to\mathbb C\)，使
\[
 \varphi_{\mu_n}(t)\longrightarrow\psi(t)
 \quad\text{对每个 }t\in\mathbb R,
\]
且 \(\psi\) 在 \(0\) 处连续，则存在唯一的 Borel 概率测度 \(\mu\)，
满足 \(\varphi_\mu=\psi\)，并且 \(\mu_n\Rightarrow\mu\)。
反过来，若 \(\mu_n\Rightarrow\mu\)，其中 \(\mu\) 是实轴上的
Borel 概率测度，则
\(\varphi_{\mu_n}(t)\to\varphi_\mu(t)\) 对每个实数 \(t\) 成立，
且极限特征函数在 \(0\) 处连续。
\end{theorem}
\begin{proof}
先证正向结论。\cref{b1:lem:30-characteristic-tightness}说明测度列为紧族。
由于 \(\mathbb R\) 是波兰空间，Prokhorov 定理给出一个弱收敛子列
\(\mu_{n_k}\Rightarrow\mu\)，其极限仍是 Borel 概率测度。
用有界连续测试函数 \(\cos(tx)\)、\(\sin(tx)\) 分别取极限，得到
\[
 \varphi_\mu(t)=\lim_{k\to\infty}\varphi_{\mu_{n_k}}(t)=\psi(t)
 \quad(t\in\mathbb R).
\]
这证明了具有所求特征函数的概率测度确实存在。
唯一性由\cref{b1:thm:30-characteristic-uniqueness}给出；
再用\cref{b1:prop:30-tight-characteristic-limit}，便得到整列
\(\mu_n\Rightarrow\mu\)。

反向结论直接用同样的余弦、正弦测试函数。
所得极限是 \(\varphi_\mu\)，它由
\cref{b1:prop:30-characteristic-basic}在整个实轴上一致连续，
特别在 \(0\) 处连续。
\end{proof}

特别地，当候选概率测度 \(\mu\) 已经给定时，
\(\mu_n\Rightarrow\mu\) 等价于所有参数 \(t\) 下的特征函数收敛。
此时无需另行检查紧性，因为极限 \(\varphi_\mu\) 在零点连续，
已经通过尾概率估计保证紧性。

\ProseParagraphBreak

零点连续性不能删去。
取 \(\mu_n=N(0,n^2)\)，则
\[
 \varphi_{\mu_n}(t)=e^{-n^2t^2/2}
 \longrightarrow
 \begin{cases}
  1,&t=0,\\
  0,&t\ne0.
 \end{cases}
\]
极限在零点不连续，因而不是任何概率测度的特征函数。
从密度也能直接看见质量逃逸：对每个固定 \(R>0\)，
\[
 \mu_n([-R,R])
 \leq\frac{2R}{\sqrt{2\pi}\,n}\longrightarrow0.
\]
任何固定紧集都装不下这列分布的大部分质量。
连续性定理中的零点条件正是把这种情形排除的条件。
下一节将计算标准化独立和的特征函数极限；
一旦极限被认作标准正态特征函数，本节的定理便把函数极限变成分布极限。
